\worksheettitle
  {Единицы длины}
  {\modulemark{M03} Учебный модуль \textbullet\ полный маршрут}

\studentnameblock

\begin{notebox}[Чему вы научитесь]
  Переводить длину в нужные единицы, не изменяя саму величину, и объяснять
  каждое преобразование через соотношение единиц.
\end{notebox}

\begin{checkpointbox}[Вход: проверьте опору]
  Заполните:
  \[
    1\text{ см}=\answerblank[15mm]\text{ мм},\qquad
    1\text{ м}=\answerblank[15mm]\text{ см},\qquad
    1\text{ км}=\answerblank[18mm]\text{ м}.
  \]
  Что меньше: сантиметр или метр? \answerblank[30mm]
\end{checkpointbox}

\begin{notebox}[Как устроен маршрут]
  Сначала разберите переход к меньшей единице, затем обратное разложение.
  В каждом задании называйте используемое соотношение. После самостоятельной
  практики исправьте ошибку и пройдите контрольную точку.
\end{notebox}

\section{Связь между единицами}

\unitsladderfigure

\begin{notebox}[Основные соотношения]
  \[
    1\text{ см}=10\text{ мм},\quad
    1\text{ дм}=10\text{ см},
  \]
  \[
    1\text{ м}=10\text{ дм}=100\text{ см},\quad
    1\text{ км}=1000\text{ м}.
  \]
  При переходе к меньшей единице число увеличивается, потому что меньших
  единиц в той же длине помещается больше.
\end{notebox}

\section{Разбираем два направления}

\begin{fillbox}[Пример 1. В одну меньшую единицу]
  Переведём $7\text{ см }9\text{ мм}$ в миллиметры.
  \[
    7\text{ см}=7\cdot10\text{ мм}=70\text{ мм},
  \]
  поэтому
  \[
    7\text{ см }9\text{ мм}=70\text{ мм}+9\text{ мм}=79\text{ мм}.
  \]
  Использовали соотношение $1\text{ см}=\answerblank[12mm]\text{ мм}$.
\end{fillbox}

\begin{fillbox}[Пример 2. Разложение по крупным единицам]
  Разложим $905\text{ см}$ на метры и сантиметры.
  В одном метре $100$ см. В числе $905$ содержится $9$ полных сотен и
  ещё $5$:
  \[
    905\text{ см}=9\text{ м }5\text{ см}.
  \]
  Проверка: $9\cdot100+5=\answerblank[22mm]\text{ см}$.
\end{fillbox}

\begin{notebox}[Алгоритм]
  \textbf{1.} Назовите исходную и требуемую единицы.
  \textbf{2.} Запишите их соотношение.
  \textbf{3.} Решите, должно ли число увеличиться или уменьшиться.
  \textbf{4.} Выполните умножение, деление или разложение.
  \textbf{5.} Запишите единицу и проверьте разумность результата.
\end{notebox}

\section{Практика с опорой}

\begin{practicebox}[1. Переведите в одну единицу]
  \begin{enumerate}[label=\alph*)]
    \item $6\text{ см }4\text{ мм}=\answerblank[25mm]\text{ мм}$;
    \item $3\text{ дм }8\text{ см}=\answerblank[25mm]\text{ см}$;
    \item $5\text{ м }8\text{ см}=\answerblank[25mm]\text{ см}$;
    \item $6\text{ км }45\text{ м}=\answerblank[28mm]\text{ м}$.
  \end{enumerate}
\end{practicebox}

\begin{practicebox}[2. Разложите]
  \begin{enumerate}[label=\alph*)]
    \item $248\text{ мм}=\answerblank[18mm]\text{ см }
      \answerblank[15mm]\text{ мм}$;
    \item $730\text{ см}=\answerblank[18mm]\text{ м }
      \answerblank[18mm]\text{ см}$;
    \item $7300\text{ м}=\answerblank[18mm]\text{ км }
      \answerblank[22mm]\text{ м}$.
  \end{enumerate}
\end{practicebox}

\section{Самостоятельная практика}

\begin{practicebox}[3. Выполните без таблицы]
  \begin{enumerate}[label=\alph*)]
    \item $4\text{ дм }7\text{ см}=\answerblank[28mm]\text{ см}$;
    \item $2\text{ м }7\text{ дм}=\answerblank[28mm]\text{ см}$;
    \item $4065\text{ м}=\answerblank[18mm]\text{ км }
      \answerblank[24mm]\text{ м}$;
    \item $1208\text{ мм}=\answerblank[18mm]\text{ м }
      \answerblank[18mm]\text{ см }
      \answerblank[15mm]\text{ мм}$.
  \end{enumerate}
\end{practicebox}

\section{Разбираем ошибку}

\begin{warningbox}[Почему получилось 645?]
  Ученик написал:
  \[
    6\text{ км }45\text{ м}=645\text{ м}.
  \]
  \textbf{Причина ошибки:} \answerblank[100mm]

  \textbf{Исправление:} \answerblank[100mm]
\end{warningbox}

\section{Контрольная точка}

\begin{checkpointbox}[4. Проверьте результат]
  \begin{enumerate}[label=\textbf{\arabic*.}]
    \item $4\text{ км }80\text{ м}=\answerblank[28mm]\text{ м}$.
    \item $605\text{ см}=\answerblank[18mm]\text{ м }
      \answerblank[18mm]\text{ см}$.
    \item Сравните $3\text{ м }8\text{ см}$ и $308\text{ см}$:
      \answerblank[45mm].
    \item Объясните, почему число при переходе от метров к сантиметрам
      увеличивается, хотя длина не меняется. \writinglines{2}
  \end{enumerate}
\end{checkpointbox}

\begin{reflectionbox}[Проверяю себя]
  \begin{itemize}[label=\choicebox]
    \item записываю соотношение единиц до вычисления;
    \item перевожу составную длину в одну единицу;
    \item раскладываю длину по крупным единицам;
    \item объясняю, почему величина не изменилась.
  \end{itemize}
\end{reflectionbox}

\begin{resultbox}[Дальнейший маршрут]
  Если преобразуете запись механически, выполните M03-R. Для тренировки
  используйте M03-H. При устойчивом результате переходите к M04 или M07.
  M03\(\star\) выполняется дополнительно.
\end{resultbox}
